% Source: https://artofproblemsolving.com/community/c3h2099887_three_more_problems
% Stable identity: aops-topic-2099887-q05
First count the numbers $0$ to $1999$, written with four digits (leading zeros do not add any $2$s). The thousands digit is $0$ or $1$. In each of the other three positions, each digit $0$--$9$ appears equally often, $\dfrac{2000}{10}=200$ times. So there are $3\cdot200=600$ twos.

Now the numbers $2000$ to $2020$: the thousands digit is $2$ in all $21$ of them; the tens digit is $2$ only in $2020$; the units digit is $2$ in $2002$ and $2012$; the hundreds digit is always $0$. That is $21+1+2=24$ twos.

In total Carla writes the digit $2$ exactly $600+24=\boxed{624}$ times.
